\documentclass{beamer}
\usepackage{graphicx}
\usepackage{xcolor}
\usepackage{tikz}

\usepackage{colortbl}     % 支持 \rowcolor
\usepackage[table]{xcolor}  % 增强颜色支持，特别是灰色背景
\usepackage{array}        % 增强表格功能（对齐方式等）
\usepackage{booktabs



% 自定义颜色
\definecolor{purpleblue}{RGB}{80, 80, 180}
\definecolor{lightgray}{RGB}{230, 230, 230}
\definecolor{novelPurple}{RGB}{98, 54, 255} 
\definecolor{darkgray}{RGB}{50, 50, 50}

% 主题设置
\usetheme{Madrid}  
\usecolortheme{dolphin} 

% 设置字体（移除 fontspec，确保 pdfLaTeX 兼容）
\setbeamerfont{title}{size=\Large,series=\bfseries}
\setbeamerfont{frametitle}{size=\LARGE,series=\bfseries}
\setbeamercolor{frametitle}{fg=white, bg=novelPurple}
\setbeamercolor{title}{fg=white, bg=purpleblue}

% 封面信息
\title[SAT Unit 3.5 ]{\textbf{SAT Unit 3.5\\ Probability}}
\author{\textbf{Jack Zeng}}
\institute{\textcolor{white}{\textbf{Novel Prep}}}
\date{\today} 

\begin{document}

% --- Title Slide ---
\begin{frame}
    \titlepage
    \vfill
    \begin{flushright}
        \textcolor{purpleblue}{\Huge \textbf{Novel Prep}}
    \end{flushright}
\end{frame}




\begin{frame}
    \begin{tikzpicture}[remember picture, overlay]
        % 设置浅色渐变背景
        \shade[top color=softPurple, bottom color=softGray] 
            (current page.north west) rectangle (current page.south east);
        
        % 添加高斯模糊光晕
        \fill[white, opacity=0.3] (current page.center) circle (4cm);
        
        % 主标题
        \node[align=center] at (current page.center) {
            {\Huge\bfseries\textcolor{purpleblue}{Novel Prep}} \\[0.5cm]
            {\Large\itshape\textcolor{lightText}{Excellence in SAT Prep}}
        };

        % 添加柔和光点（点缀）
        \foreach \x in {1,2,3,4,5} {
            \fill[purpleblue, opacity=0.2] (rand*10-5, rand*7-3.5) circle (0.1+\x*0.05);
        }

        ;
    \end{tikzpicture}
\end{frame}

\begin{frame}{Overview}
    \tableofcontents
\end{frame}


\begin{frame}
\section{Probability}
    \begin{tikzpicture}[remember picture, overlay]
        % 设置浅色渐变背景
        \shade[top color=softPurple, bottom color=softGray] 
            (current page.north west) rectangle (current page.south east);
        
        % 添加高斯模糊光晕
        \fill[white, opacity=0.3] (current page.center) circle (4cm);
        
        % 主标题
        \node[align=center] at (current page.center) {
            {\LARGE\bfseries\textcolor{purpleblue}{ Probability}} \\[0.5cm]
            {\Large\itshape\textcolor{lightText}{Excellence in SAT Prep}}
        };

        % 添加柔和光点（点缀）
        \foreach \x in {1,2,3,4,5} {
            \fill[purpleblue, opacity=0.2] (rand*10-5, rand*7-3.5) circle (0.1+\x*0.05);
        }

        ;
    \end{tikzpicture}
\end{frame}


\begin{frame}{Definition of Probability}
\small

\textbf{Probability} is a measure of how likely an event is to occur. It is defined as:

\[
P(\text{Event}) = \frac{\text{Number of favorable outcomes}}{\text{Total number of possible outcomes}}
\]

\textbf{Key Properties:}
\begin{itemize}
    \item $0 \leq P(A) \leq 1$ for any event $A$
    \item $P(\text{certain event}) = 1$
    \item $P(\text{impossible event}) = 0$
    \item $P(\text{not A}) = 1 - P(A)$
    \item For mutually exclusive events $A$ and $B$:
    \[
    P(A \text{ or } B) = P(A) + P(B)
    \]
\end{itemize}

\end{frame}

\begin{frame}{Simple Probability Examples}
\small

\textbf{Example 1:} A fair die is rolled. What is the probability of getting an even number?

\[
\text{Even numbers: } \{2, 4, 6\} \Rightarrow \text{Favorable outcomes: } 3
\]
\[
\text{Total outcomes: } 6 \Rightarrow P(\text{even}) = \frac{3}{6} = \frac{1}{2}
\]

\vspace{0.3cm}
\textbf{Example 2:} A card is drawn from a standard deck. What is the probability of drawing a heart?

\[
\text{Hearts: } 13,\quad \text{Total cards: } 52 \Rightarrow P(\text{heart}) = \frac{13}{52} = \frac{1}{4}
\]

\end{frame}


\begin{frame}{Conditional Probability (Concept)}
\small

\textbf{Conditional Probability} is the probability of event $A$ occurring \textbf{given} that event $B$ has occurred.

\[
P(A \mid B) = \frac{P(A \cap B)}{P(B)}
\]

\textbf{Key Points:}
\begin{itemize}
    \item $P(A \mid B)$ narrows the sample space to only outcomes where $B$ happens.
    \item Only defined if $P(B) \ne 0$.
\end{itemize}

\textbf{Visual:}

Think of "restricting" the universe to $B$, and checking how many of those are also $A$


\end{frame}

\begin{frame}{\textbf{Probability - Score Distribution Table}}
\small

The same 20 contestants, on each of 3 days, answered 5 questions in order to win a prize. Each contestant received 1 point for each correct answer. The number of contestants receiving a given score on each day is shown in the table below.

\begin{center}
\begin{tabular}{|c|c|c|c|c|}
\hline
\textbf{Score} & \textbf{Day 1} & \textbf{Day 2} & \textbf{Day 3} & \textbf{Total} \\
\hline
5 out of 5 & 2 & 2 & 3 & 7 \\
4 out of 5 & 3 & 3 & 3 & 9 \\
3 out of 5 & 4 & 5 & 3 & 13 \\
2 out of 5 & 6 & 5 & 4 & 15 \\
1 out of 5 & 2 & 4 & 3 & 9 \\
0 out of 5 & 3 & 1 & 2 & 6 \\
\hline
\textbf{Total} & 20 & 20 & 20 & 60 \\
\hline
\end{tabular}
\end{center}

No contestant received the same score on two different days. If a contestant is selected at random, what is the probability that the selected contestant received a score of 5 on Day 2 or Day 3, given that the contestant received a score of 5 on one of the three days?

\end{frame}

\begin{frame}{Answer: Conditional Probability}
\small

\textbf{Problem Context:}
\begin{itemize}
    \item 20 contestants each day, over 3 days (no repeat scores).
    \item Table shows number of contestants scoring 0–5 on each day.
    \item Total scores of 5 across all days: $2 + 2 + 3 = 7$
\end{itemize}


\textbf{Given:} A contestant received a score of $5$ on \textbf{one of the days}.

\textbf{Find:} Probability that this contestant scored $5$ on Day 2 or Day 3.


\textbf{Step 1: Count favorable outcomes}
\[
\text{Day 2: } 2 \qquad \text{Day 3: } 3 \qquad \text{Total favorable} = 2 + 3 = 5
\]

\vspace{0.1cm}
\textbf{Step 2: Count total outcomes}
\[
\text{Total contestants with score of 5: } 7
\]

\vspace{0.1cm}
\textbf{Answer:}
\[
P(\text{Day 2 or Day 3} \mid \text{Score 5}) = \frac{5}{7}
\]

\textbf{Final Answer:} $\boxed{\frac{5}{7}}$

\end{frame}



\begin{frame}
\section{Practices}
    \begin{tikzpicture}[remember picture, overlay]
        % 设置浅色渐变背景
        \shade[top color=softPurple, bottom color=softGray] 
            (current page.north west) rectangle (current page.south east);
        
        % 添加高斯模糊光晕
        \fill[white, opacity=0.3] (current page.center) circle (4cm);
        
        % 主标题
        \node[align=center] at (current page.center) {
            {\Huge\bfseries\textcolor{purpleblue}{Practices}} \\[0.5cm]
            {\Large\itshape\textcolor{lightText}{Excellence in SAT Prep}}
        };

        % 添加柔和光点（点缀）
        \foreach \x in {1,2,3,4,5} {
            \fill[purpleblue, opacity=0.2] (rand*10-5, rand*7-3.5) circle (0.1+\x*0.05);
        }

        ;
    \end{tikzpicture}
\end{frame}


\begin{frame}{Question}
The table summarizes the distribution of age and favorite sport for 150 members of a country club.

\begin{center}
\begin{tabular}{|c|c|c|c|c|}
\hline
 & Squash & Tennis & Golf & Total \\
\hline
30–44 years & 18 & 10 & 22 & 50 \\
45–59 years & 11 & 14 & 25 & 50 \\
60+ years & 21 & 26 & 3 & 50 \\
\hline
Total & 50 & 50 & 50 & 150 \\
\hline
\end{tabular}
\end{center}

One of these members will be selected at random. What is the probability of selecting a member whose favorite sport is tennis, given that the member is at least 45 years of age? (Express your answer as a decimal or fraction, not as a percent.)
\end{frame}
\begin{frame}{Answer}
\textbf{Step 1. Identify the relevant group:}
\[
\text{At least 45 years of age} = (45–59) + (60+) = 50 + 50 = 100.
\]

\textbf{Step 2. Find the total number of tennis players in that group:}
\[
14 (\text{45–59}) + 26 (\text{60+}) = 40.
\]

\textbf{Step 3. Calculate the probability:}
\[
P(\text{Tennis} \mid \text{At least 45}) = \frac{40}{100} = 0.4.
\]

\textbf{Answer:} $\boxed{0.4}$.
\end{frame}


\begin{frame}{Question}
A new medical test is designed to give a positive indicator to a patient who has a certain virus and a negative indicator to a patient who does not. The table summarizes the results for 1,000 patients who have taken this test.

\begin{center}
\begin{tabular}{|c|c|c|c|}
\hline
 & Negative Indicator & Positive Indicator & Total \\
\hline
Has virus & 30 & 370 & 400 \\
Does not have virus & 550 & 50 & 600 \\
\hline
Total & 580 & 420 & 1,000 \\
\hline
\end{tabular}
\end{center}

Based on the table, what is the probability that the test gives the wrong indicator?

\begin{enumerate}[A.]
    \item 5\%
    \item 8\%
    \item 10\%
    \item 12\%
\end{enumerate}
\end{frame}


\begin{frame}{Answer}
\textbf{Step 1. Count the number of wrong indicators:}
\begin{itemize}
    \item False negatives (Has virus, but test negative): 30
    \item False positives (Does not have virus, but test positive): 50
\end{itemize}
\[
30 + 50 = 80
\]

\textbf{Step 2. Divide by the total number of patients:}
\[
\frac{80}{1,000} = 0.08
\]

\textbf{Step 3. Answer:}
\[
0.08 = 8\%
\]

\textbf{Answer:} B.
\end{frame}


\begin{frame}{Question}
\small
The table shows the results of a survey in which students at three colleges were asked whether they live on campus or off campus.

\begin{center}
\begin{tabular}{|c|c|c|c|}
\hline
 & Off campus & On campus & Total \\
\hline
College A & 335 & 689 & 1,024 \\
College B & 124 & 482 & 606 \\
College C & 596 & 274 & 870 \\
\hline
Total & 1,055 & 1,445 & 2,500 \\
\hline
\end{tabular}
\end{center}

If one of the students surveyed is selected at random, which of the following probabilities has the greatest value?

\begin{enumerate}[A.]
    \item The probability that the student lives off campus, given that the student attends College C.
    \item The probability that the student lives on campus, given that the student attends College C.
    \item The probability that the student attends College C, given that the student lives off campus.
    \item The probability that the student attends College C, given that the student lives on campus.
\end{enumerate}
\end{frame}


\begin{frame}{Answer}
\textbf{Step 1. Calculate each conditional probability:}

\begin{itemize}
    \item A: $\frac{596}{870} \approx 0.685$
    \item B: $\frac{274}{870} \approx 0.315$
    \item C: $\frac{596}{1055} \approx 0.565$
    \item D: $\frac{274}{1445} \approx 0.19$
\end{itemize}

\textbf{Step 2. Compare:}
\[
0.685 > 0.565 > 0.315 > 0.19
\]

\textbf{Answer:} A.
\end{frame}








\begin{frame}{}
    \begin{tikzpicture}[remember picture, overlay]
        % 背景渐变色：蓝紫到白
        \shade[top color=novelPurple!40, bottom color=white]
            (current page.north west) rectangle (current page.south east);

        % 中心柔光
        \fill[white, opacity=0.15] (current page.center) circle (5cm);

        % 柔和光点点缀
        \foreach \x in {1,2,3,4,5} {
            \fill[novelPurple, opacity=0.12] (rand*10-5, rand*7-3.5) circle (0.1+\x*0.05);
        }
    \end{tikzpicture}

    \centering
    \vspace{5em}

    {\Huge \textbf{\textcolor{novelPurple}{Thank You!}}} \\[1em]
    {\large \textcolor{gray}{We appreciate your time and focus.}} \\[3em]

    {\LARGE \textbf{\textcolor{novelPurple}{\textsf{Novel Prep}}}}

    \vfill
\end{frame}


\end{document}
